\ifdefined\gkver\else\def\gkver{student}\fi
\documentclass[\gkver]{gaokaozhenti}
\begin{document}

\chapter{2018年上海}
%% recalled: true

\begin{enumerate}
\item
%% number: 1
%% paperName: 2018年上海等级考第 1 题 3 分
%% typeId: 1
%% type: 单选题
%% chapter: 原子和原子核
%% point: 天然放射性
%% method: 无
%% score: 3
%% degree: 659
%% duplicateId: 0
%% body:
$\alpha$ 粒子是\xzanswer{C}
\fourchoices[answer=C]
{分子}
{原子}
{原子核}
{光子}
%% answer: C
%% memo:
\memoanswer{
本题原卷及材料中未提供详解，待补充。
}

\item
%% number: 2
%% paperName: 2018年上海等级考第 2 题 3 分
%% typeId: 1
%% type: 单选题
%% chapter: 光学
%% point: 光电效应
%% method: 无
%% score: 3
%% degree: 870
%% duplicateId: 0
%% body:
能解释光电效应的学说是\xzanswer{D}
\fourchoices[answer=D]
{牛顿的微粒说}
{惠更斯的光的波动说}
{麦克斯韦的光的电磁说}
{爱因斯坦的光子说}
%% answer: D
%% memo:
\memoanswer{
本题原卷及材料中未提供详解，待补充。
}

\item
%% number: 3
%% paperName: 2018年上海等级考第 3 题 3 分
%% typeId: 1
%% type: 单选题
%% chapter: 原子和原子核
%% point: 人工核反应
%% method: 无
%% score: 3
%% degree: 970
%% duplicateId: 0
%% body:
核反应方程 $\ce{^{4}_{2}He + ^{9}_{4}Be -> ^{12}_{6}C + ^{A}_{Z}X}$ 中的 $\ce{^{A}_{Z}X}$ 是\xzanswer{D}
\fourchoices[answer=D]
{$\ce{^{1}_{1}H}$}
{$\ce{^{0}_{1}e}$}
{$\ce{^{0}_{-1}e}$}
{$\ce{^{1}_{0}n}$}
%% answer: D
%% memo:
\memoanswer{
本题原卷及材料中未提供详解，待补充。
}

\item
%% number: 4
%% paperName: 2018年上海等级考第 4 题 3 分
%% typeId: 1
%% type: 单选题
%% chapter: 机械波
%% point: 波长频率波速
%% method: 无
%% score: 3
%% degree: 820
%% duplicateId: 0
%% body:
如图所示，手抓住绳的一端上下振动，产生沿绳传播的机械波，增大手的振动频率，该波的\xzanswer{C}
\onepicture[h=1.5cm, align=right]{figs/04.png}
\fourchoices[answer=C]
{波速增大}
{波长增大}
{波速不变}
{波长不变}
%% answer: C
%% memo:
\memoanswer{
本题原卷及材料中未提供详解，待补充。
}

\item
%% number: 5
%% paperName: 2018年上海等级考第 5 题 3 分
%% typeId: 1
%% type: 单选题
%% chapter: 磁场
%% point: 安培力
%% method: 无
%% score: 3
%% degree: 930
%% duplicateId: 0
%% body:
算式“$3\UT\cdot4\UA\cdot2\Um$”的计算结果是\xzanswer{A}
\fourchoices[answer=A]
{24 N}
{24 V}
{24 W}
{24 J}
%% answer: A
%% memo:
\memoanswer{
本题原卷及材料中未提供详解，待补充。
}

\item
%% number: 6
%% paperName: 2018年上海等级考第 6 题 3 分
%% typeId: 1
%% type: 单选题
%% chapter: 气体性质
%% point: 物体的内能
%% method: 无
%% score: 3
%% degree: 701
%% duplicateId: 0
%% body:
在物理学中，有些物理量与物体的状态有关，另一些物理量则与物理过程有关，下列物理量中与物理过程有关的是\xzanswer{A}
\fourchoices[answer=A]
{热量}
{内能}
{压强}
{体积}
%% answer: A
%% memo:
\memoanswer{
本题原卷及材料中未提供详解，待补充。
}

\item
%% number: 7
%% paperName: 2018年上海等级考第 7 题 3 分
%% typeId: 1
%% type: 单选题
%% chapter: 电磁感应
%% point: 楞次定律
%% method: 无
%% score: 3
%% degree: 810
%% duplicateId: 0
%% body:
如图所示，线圈 L、M、R 平行共轴放置，M 中通有电流强度为 $I$ 的稳恒电流。若 M 按箭头方向向左运动，从左侧看到的 L、R 两线圈中的感应电流方向为\xzanswer{A}
\onepicture[h=3.0cm]{\includesvg{figs/07}}
\fourchoices[answer=A]
{L 顺时针，R 逆时针}
{L 逆时针，R 顺时针}
{L、R 都是顺时针}
{L、R 都是逆时针}
%% answer: A
%% memo:
\memoanswer{
本题原卷及材料中未提供详解，待补充。
}

\item
%% number: 8
%% paperName: 2018年上海等级考第 8 题 3 分
%% typeId: 1
%% type: 单选题
%% chapter: 曲线运动
%% point: 人造地球卫星
%% method: 无
%% score: 3
%% degree: 860
%% duplicateId: 0
%% body:
行星绕恒星做匀速圆周运动的线速度大小取决于\xzanswer{B}
\fourchoices[answer=B]
{行星轨道半径}
{恒星质量及行星轨道半径}
{行星及恒星质量}
{行星质量及轨道半径}
%% answer: B
%% memo:
\memoanswer{
本题原卷及材料中未提供详解，待补充。
}

\item
%% number: 9
%% paperName: 2018年上海等级考第 9 题 4 分
%% typeId: 1
%% type: 单选题
%% chapter: 力物体的平衡
%% point: 三力平衡
%% method: 无
%% score: 4
%% degree: 730
%% duplicateId: 0
%% body:
三个大小都为 $F$ 的力作用于一质点，其中两力始终相互垂直，另一个力方向可变，该质点\xzanswer{A}
\onepicture[h=3.0cm]{\includesvg{figs/09}}
\fourchoices[answer=A]
{不能平衡}
{能否平衡取决于 $\theta$ 的大小}
{能否平衡取决于 $F$ 的大小}
{能否平衡取决于 $F$ 和 $\theta$ 的大小}
%% answer: A
%% memo:
\memoanswer{
本题原卷及材料中未提供详解，待补充。
}

\item
%% number: 10
%% paperName: 2018年上海等级考第 10 题 4 分
%% typeId: 1
%% type: 单选题
%% chapter: 电场
%% point: 电场中的图象
%% method: 无
%% score: 4
%% degree: 870
%% duplicateId: 0
%% body:
在 $x$ 轴上电场强度 $E$ 随位置变化的情况如图所示，$E>0$ 表示电场方向与 $x$ 轴正方向一致。一正电荷由 $x_{0}$ 出发向 $x$ 轴正方向运动到 $x_{1}$ 的过程中，其电势能\xzanswer{C}
\onepicture[h=2.8cm]{\includesvg{figs/10}}
\fourchoices[answer=C]
{始终增大}
{始终减小}
{先减小后增大}
{先增大后减小}
%% answer: C
%% memo:
\memoanswer{
本题原卷及材料中未提供详解，待补充。
}

\item
%% number: 11
%% paperName: 2018年上海等级考第 11 题 4 分
%% typeId: 1
%% type: 单选题
%% chapter: 功和能
%% point: 动能势能
%% method: 无
%% score: 4
%% degree: 763
%% duplicateId: 0
%% body:
撑杆运动员在最高点上放开撑杆，并水平越过横杆，若撑杆的弹性势能为 $E_{e}$，运动员的动能为 $E_{k}$，运动员的重力势能为 $E_{p}$，则在运动员放手瞬间，三种能量的相对大小为\xzanswer{B}
\fourchoices[answer=B, ispicture=true, h=2.6cm]
{\includesvg{figs/11a}}
{\includesvg{figs/11b}}
{\includesvg{figs/11c}}
{\includesvg{figs/11d}}
%% answer: B
%% memo:
\memoanswer{
本题原卷及材料中未提供详解，待补充。
}

\item
%% number: 12
%% paperName: 2018年上海等级考第 12 题 4 分
%% typeId: 1
%% type: 单选题
%% chapter: 电路
%% point: 动态分析
%% method: 无
%% score: 4
%% degree: 670
%% duplicateId: 0
%% body:
恒压电源的内阻可视为零，普通电源的内阻不可忽略。为判别一电源是哪种电源，某同学连接了如图所示实验电路，并断开电键，设想通过观察电键闭合后的电流表作出判断。此方案\xzanswer{D}
\onepicture[h=3.4cm, align=right]{\includesvg{figs/12}}
\fourchoices[answer=D]
{无效。因为无论是哪种电路，电流表示数都变大}
{无效。因为无论是哪种电路，电流表示数都变小}
{有效。若是恒压电源，电流表示数不变，若是普通电源，则电流表示数变大}
{有效。若是恒压电源，电流表示数不变，若是普通电源，则电流表示数变小}
%% answer: D
%% memo:
\memoanswer{
本题原卷及材料中未提供详解，待补充。
}

\item
%% number: 13
%% paperName: 2018年上海等级考第 13 题 4 分
%% typeId: 3
%% type: 填空题
%% chapter: 光学
%% point: 光的干涉
%% method: 无
%% score: 4
%% degree: 901
%% duplicateId: 0
%% body:
分别用红光和紫光照射同一双缝，保持其他条件不变，得到的干涉条纹间距较宽的是\tkanswer[1.0cm]{红}光，当改用白光照射时，除中央亮条纹外，两侧的亮条纹是\tkanswer[1.0cm]{彩}色的。
%% answer: 红，彩
%% memo:
\memoanswer{
本题原卷及材料中未提供详解，待补充。
}

\item
%% number: 14
%% paperName: 2018年上海等级考第 14 题 4 分
%% typeId: 3
%% type: 填空题
%% chapter: 气体性质
%% point: 分子动理论
%% method: 无
%% score: 4
%% degree: 880
%% duplicateId: 0
%% body:
伽耳顿板可以演示统计规律，若让一个小钢珠从漏斗形入口落下，与铁钉发生一系列碰撞，最终会落在哪个槽中是\tkanswer[1.2cm]{不能}（选填“能”或“不能”）事先确定的；若让大量小钢珠从漏斗形入口落下，落在各个狭槽中的小钢珠数目\tkanswer[2.4cm]{中间较多}（选填“基本相等”、“中间较多”或“两侧较多”）。
\onepicture[h=3.4cm, align=right]{\includesvg{figs/14}}
%% answer: 不能，中间较多
%% memo:
\memoanswer{
本题原卷及材料中未提供详解，待补充。
}

\item
%% number: 15
%% paperName: 2018年上海等级考第 15 题 4 分
%% typeId: 3
%% type: 填空题
%% chapter: 直线运动
%% point: 运动图象
%% method: 无
%% score: 4
%% degree: 870
%% duplicateId: 0
%% body:
一质点从静止开始沿直线运动的速度与时间的关系如图所示，该质点远离起点的时间范围为\tkanswer[1.8cm]{$0\sim7$}$\Us$，质点在 $0\sim9\Us$ 内的位移为\tkanswer[1.2cm]{14}$\Um$。
\onepicture[h=2.8cm, align=right]{\includesvg{figs/15}}
%% answer: $0\sim7$，14
%% memo:
\memoanswer{
本题原卷及材料中未提供详解，待补充。
}

\item
%% number: 16
%% paperName: 2018年上海等级考第 16 题 4 分
%% typeId: 3
%% type: 填空题
%% chapter: 电路
%% point: 闭合电路欧姆定律
%% method: 无
%% score: 4
%% degree: 870
%% duplicateId: 0
%% body:
将电源与一定值电阻 A 相连，外电压为 $U$，若改用阻值为其 3 倍的定值电阻 B，外电压变为 $2U$，则此时电流为原来的\tkanswer[1.2cm]{$\frac{2}{3}$}倍，电源电动势为\tkanswer[1.2cm]{$4U$}。
%% answer: $\frac{2}{3}$，$4U$
%% memo:
\memoanswer{
本题原卷及材料中未提供详解，待补充。
}

\item
%% number: 17
%% paperName: 2018年上海等级考第 17 题 4 分
%% typeId: 3
%% type: 填空题
%% chapter: 气体性质
%% point: 查理定律
%% method: 无
%% score: 4
%% degree: 634
%% duplicateId: 0
%% body:
利用如图所示装置可研究一定质量的气体体积不变时温度与压强的关系，玻璃管 A、B 及橡皮弯管内装有水银，烧瓶内封闭一定质量的气体。当温度改变时，需上下移动 A 管，使 B 管液面\tkanswer[3.0cm]{回到初始位置}，已知大气压强为 $76\UcmHg$，还需要使用的测量工具是\tkanswer[1.5cm]{刻度尺}和\tkanswer[1.5cm]{温度计}。
\onepicture[h=3.4cm, align=right]{\includesvg{figs/17}}
%% answer: 回到初始位置，刻度尺，温度计
%% memo:
\memoanswer{
本题原卷及材料中未提供详解，待补充。
}

\item
%% number: 18
%% paperName: 2018年上海等级考第 18 题 12 分
%% typeId: 4
%% type: 实验题
%% chapter: 运动和力
%% point: 验证牛顿第二定律
%% method: 无
%% score: 12
%% degree: 701
%% duplicateId: 0
%% body:
用“DIS 研究加速度和质量关系”的实验装置如图所示。
\begin{enumerate}
	\item
	图中使用的是\tkanswer[1.8cm]{位移}传感器；
	\item
	安装实验装置并运行 DIS 软件，进入本实验界面；随后的操作顺序是\tkanswer[1.6cm]{CAB}（用字母排列）：

	A．点击“选择区域”，在图线中选择适当区域，计算机得出该区域的加速度，并记入表格；在表格中输入小车总质量；

	B．保持钩码个数不变，增加配重片，重复实验得到几组数据；

	C．点击“开始记录”并释放小车，完成运动后点击“停止记录”，得到 $v-t$ 图像。
	\item
	某同学不使用图中的传感器，而是在轨道 A 处固定一光电门传感器，实验时改变配重片数目，记录小车总质量 $m$，由静止释放小车，测出小车过 A 时的瞬时速度 $v$，由实验数据得到 $v^{2}-\frac{1}{m}$ 图，验证加速度与质量的关系。实验中每次释放小车的位置\tkanswer[1.5cm]{必须相同}（选填“必须相同”或“可以不同”），原因是\tkanswer[4.6cm]{由 $v^{2}=2as$ 可知，只有在 $s$ 相同的情况下，才有 $v^{2}$ 正比于 $a$，使得 $v^{2}-\frac{1}{m}$ 图也是过原点一条直线，从而证明结论}。
\end{enumerate}
\onepicture[w=9cm, align=right]{\includesvg{figs/18}}
%% answer: （1）位移（2）CAB（3）必须相同，由 $v^{2}=2as$ 可知，只有在 $s$ 相同的情况下，才有 $v^{2}$ 正比于 $a$，使得 $v^{2}-\frac{1}{m}$ 图也是过原点一条直线，从而证明结论。
\jdanswer{
\begin{enumerate}
	\item
	位移

	\item
	CAB

	\item
	必须相同，由 $v^{2}=2as$ 可知，只有在 $s$ 相同的情况下，才有 $v^{2}$ 正比于 $a$，使得 $v^{2}-\frac{1}{m}$ 图也是过原点一条直线，从而证明结论。

\end{enumerate}
}
%% memo:
\memoanswer{
本题原卷及材料中未提供详解，待补充。
}

\item
%% number: 19
%% paperName: 2018年上海等级考第 19 题 12 分
%% typeId: 6
%% type: 计算题
%% chapter: 电磁感应
%% point: 电磁感应中的匀变速
%% method: 无
%% score: 12
%% degree: 859
%% duplicateId: 0
%% body:
如图，匀强磁场的磁感应强度为 $B$，方向与水平面垂直。边长为 $L$ 的正方形导线框位于水平面内，$t=0$ 时，线框从磁场左边缘由静止开始向右做匀加速运动；$t=T$ 时线框刚要完全进入磁场，此时线框中感应电流强度为 $I_{0}$。求：
\begin{enumerate}
	\item
	线框做匀加速运动的加速度大小 $a$；
	\item
	线框的电阻 $R$；
	\item
	线框进入磁场过程中受到的安培力随时间变化的关系式 $F_{A}(t)$。
\end{enumerate}
\onepicture[h=2.6cm, align=right]{\includesvg{figs/19}}
%% answer: （1）$a=\frac{2L}{T^{2}}$（2）$R=\frac{2BL^{2}}{TI_{0}}$（3）$F_{A}(t)=\frac{BLI_{0}}{T}t$
\jdanswer{
\begin{enumerate}
	\item
	$a=\frac{2L}{T^{2}}$

	\item
	$R=\frac{2BL^{2}}{TI_{0}}$

	\item
	$F_{A}(t)=\frac{BLI_{0}}{T}t$

\end{enumerate}
}
%% memo:
\memoanswer{
（1）由运动学公式：$L=\frac{1}{2}aT^{2}$

可得：$a=\frac{2L}{T^{2}}$

（2）感应电流的大小为 $I_{0}=\frac{BLv}{R}$

由运动学公式 $v=aT=\frac{2L}{T}$

由以上两式可得：$R=\frac{BLv}{I_{0}}=\frac{2BL^{2}}{TI_{0}}$

（3）$F_{A}=BIL=B\frac{BLv}{R}L=\frac{B^{2}L^{2}a}{R}t$

将（2）中求得的 $a$、$R$ 表达式代入上式，可得

$F_{A}(t)=\frac{BLI_{0}}{T}t$
}

\item
%% number: 20
%% paperName: 2018年上海等级考第 20 题 14 分
%% typeId: 6
%% type: 计算题
%% chapter: 功和能
%% point: 交通工具启动
%% method: 无
%% score: 14
%% degree: 630
%% duplicateId: 0
%% body:
某汽车匀速行驶时发动机和传动与变速系统内的功率分配关系如图所示。图中数据为车以 $v_{0}=72\Ukmh$ 的速率匀速行驶时的功率。汽车行驶时所受空气阻力与瞬时速率的关系为 $f_{a}=kv^{2}$（$k$ 为恒量），所受路面的阻力 $f_{s}$ 大小恒定。求：
\begin{enumerate}
	\item
	恒量 $k$ 的单位；
	\item
	汽车以 $v_{0}$ 匀速运动时，发动机的输出功率 $P_{0}$；
	\item
	汽车以 $v_{0}$ 匀速运动时受到的驱动力 $F_{0}$ 的大小；
	\item
	若汽车发动机最大输出功率 $P_{\max}=150\UkW$，水泵功率 $P_{1}$ 恒定，传动与变速系统因内部机件摩擦而损耗的功率 $P_{2}$ 与汽车的行驶速率成正比。通过计算说明该汽车能否以速率 $3v_{0}$ 匀速行驶。
\end{enumerate}
\onepicture[w=10cm, align=right]{\includesvg{figs/20}}
%% answer: （1）kg/m（2）17 kW（3）500 N（4）此时所需的功率为 165 kW > 150 kW，因此不能达到 $3v_{0}$。
\jdanswer{
\begin{enumerate}
	\item
	$\Ukgm$

	\item
	$17\UkW$

	\item
	$500\UN$

	\item
	此时所需的功率为 $165\UkW>150\UkW$，因此不能达到 $3v_{0}$。

\end{enumerate}
}
%% memo:
\memoanswer{
（1）由 $f_{a}=kv^{2}$ 可知，$k$ 的单位为 $\UNsqmq$，也可以写成 $\Ukgm$。

（2）由图可知 $P_{\text{出}}=(69-52)\UkW=17\UkW$

（3）$F=\frac{P_{3}+P_{4}}{v_{0}}=\frac{5000+5000}{20}\UN=500\UN$

（4）由于速度变为原来的 3 倍，且机件摩擦和地面阻力的损耗功率与速度成正比，可知

$P_{2}'=3P_{2}=12\UkW$，$P_{4}'=3P_{4}=15\UkW$

由 $f_{a}=kv^{2}$ 可知 $P_{3}'=9f_{a}\cdot3v_{0}=27P_{3}=135\UkW$

$P_{\text{出}}'=P_{1}+P_{2}'+P_{3}'+P_{4}'=(3+12+135+15)\UkW=165\UkW>150\UkW$

因此汽车无法达到 $3v_{0}$ 的速度。
}

\end{enumerate}

\end{document}
